\section{国际融资法律制度}
\begin{enumerate}
  \item 特殊的国际贷款方式。
        \begin{enumerate}
          \item 政府贷款。援助性质，长期低息甚至无息。不具有可考性。
          \item 国际金融机构贷款。国际货币基金组织、世界银行。一般不考。个案不太具有可考性。
          \item 国际商业贷款。
                \begin{enumerate}
                  \item 国际银团贷款。统一的贷款条件。巨额长期贷款。只考过 1 次。
                        \begin{enumerate}
                          \item 直接式银团贷款。没有牵头银行，银团和借款人直接签订。
                          \item 间接式银团款。借款人 \yjt 牵头银行 \yjt 银团。
                          \item 各银行成员之间责任是按份还是连带看合同。
                        \end{enumerate}
                  \item 国际定期贷款。
                  \item 国际项目贷款/融资。
                        \begin{enumerate}
                          \item 对某一特定的工程项目发放的贷款，以项目建成后的经济收益还本付息。
                          \item 周期长、风险最大、利率最高。
                          \item 可组合。国际银团项目贷款。
                        \end{enumerate}
                \end{enumerate}
          \item 特别提款权。偿还贷款、政府间结算、国际储备。反正就是不能用于支付。
        \end{enumerate}
  \item 国际融资担保方式。重点的重点。
        \begin{kk}{见索即付保函/独立保函}
          \begin{enumerate}
            \item 3 方当事人。申请人、开立人、受益人。
            \item 要立马想到 3 性。
                  \begin{enumerate}
                    \item 独立性。效力独立于基础合同/基础交易。认为应当适用担保法一般保证和连带保证责任的，不予支持。
                    \item 连带性。连带保证，不能行使先诉抗辩权。
                    \item 支付无条件性。受益人凭违约证明 + 索偿通知书找开立人。单函、单单表面一致，开立人即向受益人付钱。
                          表面不完全一致，但并不导致歧义，应认定构成表面相符 \dui。
                  \end{enumerate}
            \item vs 信用证。是一种支付方式，见索即付保函是一种担保方式。
            \item 缺点。同信用证。
                  \begin{enumerate}
                    \item 偿付性很强，开立人不能进行实质审查。
                    \item 容易出现保函欺诈 \yjt 反欺诈手段。止付令。
                    \item 管辖。自治 \yjt 被告住所地、开立人住所地法院。
                    \item 法律适用。涉外可以协议选 \yjt 开立人经居地。分支机构开立的 \yjt 分支机构登记地法。
                  \end{enumerate}
          \end{enumerate}
        \end{kk}
        \begin{enumerate}
          \item 备用信用证。本质是一个见索即付保函。
                \begin{enumerate}
                  \item 若债务人履行债务，本信用证备而不用。
                  \item 开立人只能是银行，是独立连带的保证。
                \end{enumerate}
          \item 意愿书。一种担保文件，小名安慰信，没有法律约束力，有道义上的约束力。
          \item 浮动抵押。和民法所学的浮动抵押一样。
        \end{enumerate}
\end{enumerate}
